\paragraph{~\trDeligneMumfordStasheff{}\trCompactification（又称\trAssociahedron）。}\label{sub:dm-stasheff}
作为集合，$\scr{R}^{d+1}$ 的\trDeligneMumfordStasheff{}\trCompactification{}是
\begin{equation}\label{eq:dm-compactification}
\bar{\scr{R}}^{d+1}=\coprod_T\scr{R}^T.
\end{equation}
这里，求并取遍所有稳定的 $d$ \trLeafedTree{}$T$，并且约定
$\scr{R}^T=\prod_{v\in\mathrm{Ve}(T)}\scr{R}^{|v|}$。为了给\eqref{eq:dm-compactification} 配备\trTopology{}与\trSmoothStructure，我们作\trStriplikeEnd{}的某个\trUniversal{}选取，并使用\trGluingConstruction。也就是说，取一棵稳定\trTree{}$T$，并对每个\trVertex{}$v$ 配上\trUniversal\trFamily{}$\scr{S}^{|v|}\to\scr{R}^{|v|}$。把这些\trGluing{}起来，便得到一个 $(d+1)$ \trPointedDisc{}的\trFamily
\begin{equation}\label{eq:gluing-family}
\scr{S}\longrightarrow\scr{R}=(-1;0)^{\mathrm{Ed}^{\mathrm{int}}(T)}\times\scr{R}^T,
\end{equation}
按定义，它被一个\trSmooth{}的\textbf{\trGluingMap}（gluing map）分类：
\begin{equation}\label{eq:gluing-map}
\gamma^T:\scr{R}\longrightarrow\scr{R}^{d+1}.
\end{equation}
如前所述，可以容许（部分或全部）\trGluingParameter{}变为零。结果是一个\trDegenerate{}的\trGluing\trSurface，它是由\trCollapse{}后的\trTree{}$\bar{T}$ 所支配的\trPointedDisc{}的\trDisjointUnion。显然，这确定了\eqref{eq:dm-compactification} 的 $\bar{T}$ \trStratum{}中的一个点，从而给出诸\trGluingMap{}的一个\trCanonical{}扩张
\begin{equation}\label{eq:extended-gluing-map}
\bar\gamma^T:\bar{\scr{R}}=(-1;0]^{\mathrm{Ed}^{\mathrm{int}}(T)}\times\scr{R}^T\longrightarrow\bar{\scr{R}}^{d+1}.
\end{equation}
在子集 $\{0\}^{\mathrm{Ed}^{\mathrm{int}}(T)}\times\scr{R}^T$ 上，这代表的是平凡的（最\trDegenerate{}的）\trGluing{}运算，因此在\eqref{eq:dm-compactification} 的 $T$ \trStratum{}上恰恰给出恒同。\trDeligneMumfordStasheff{}\trCompactification{}的\trTopology{}，就是由对所有 $T$ 的映射\eqref{eq:extended-gluing-map} 诱导出来的那个：也就是说，从 $\bar{\scr{R}}^{d+1}$ 到任意空间的映射为连续，当且仅当它与诸 $\bar\gamma^T$ 的全部\trComposition{}都连续。可微结构由类似的性质刻画，只需把“连续”换成“\trSmooth”。这使得 $\bar{\scr{R}}^{d+1}$ 成为一个\trCompact{}的\trCorneredBoundary\trManifold，而且这一结构不依赖于\trStriplikeEnd{}的\trUniversal{}选取；此外，\eqref{eq:dm-compactification} 正是内部与\trConnected\trBoundaryStratum{}的分解。
我们刚刚陈述的这些性质，从\trGluingConstruction{}本身来看并不是一目了然的（事实上，就连 $\bar{\scr{R}}^{d+1}$ 上的\trTopology{}是否\trHausdorff，可能都非显然）。验证它们最好的办法，是从 $\bar{\scr{R}}^{d+1}$ 作为复\trDeligne--\trMumford{}空间 $\bar{\scr{M}}_{0,d+1}$ 的\trRealLocus{}$\bar{\scr{M}}_{0,d+1}(\RR)$ 的一部分这一替代定义出发；该定义在 [58]、[60] 中有解释。以这种方式引入时，$\bar{\scr{R}}^{d+1}$ 便自带\trSmooth\trCompact\trCorneredBoundary\trManifold{}的\trCanonical{}结构。
\begin{lemma}\label{thm:dm-smooth-structure}
为 $\bar{\scr{R}}^{d+1}$ 配备由其到 $\bar{\scr{M}}_{0,d+1}(\RR)$ 的\trEmbedding{}所带来的\trSmoothStructure。那么诸扩张\trGluingMap{}$\bar\gamma^T$ 都是\trSmooth{}的，并且在所有\trGluingParameter{}皆为零的子集附近都是\trLocalDiffeomorphism。
\end{lemma}
\begin{proof}[证明概要]
首先回顾，\trEmbedding{}$\scr{R}^{d+1}\subset\bar{\scr{M}}_{0,d+1}(\RR)\subset\bar{\scr{M}}_{0,d+1}$ 是通过把一个\trPointedDisc{}“加倍”为一个带\trAntiholomorphic{}\trInvolution{}的\trPunctured{}的 $\CP^1$ 而实现的。\trDisk{}上的诸\trStriplikeEnd{}同样可以加倍，从而给出 $\CP^1$ 上诸\trPunctured{}点附近的局部\trHolomorphic{}坐标。将诸\trGluingParameter{}\trComplexification，便可把\eqref{eq:extended-gluing-map} 扩张为映射
\begin{equation}\label{eq:complexified-gluing-map}
(\bar\gamma^T)^{\CC}:\mathrm{int}(D)^{\mathrm{Ed}^{\mathrm{int}}(T)}\times\scr{R}^T\longrightarrow\bar{\scr{M}}_{0,d+1},
\end{equation}
其中 $\mathrm{int}(D)=D\setminus\partial D$ 是开的单位\trDisk。更确切地说，若取定 $\scr{R}^T$ 上的一点，则\trComplexification{}的\trGluingConstruction{}便在 $\mathrm{int}(D)^{\mathrm{Ed}^{\mathrm{int}}(T)}$ 之上产出一个由稳定的\trPointedDisc{}组成的\trHolomorphic\trFamily。\eqref{eq:complexified-gluing-map} 将是这个\trFamily{}的\trClassifyingMap；因此它关于诸\trGluingParameter{}是\trHolomorphic{}的，并且在所有\trGluingParameter{}皆非零的开子集上是\trSmooth{}的。要得出它在处处都\trSmooth{}的结论，可利用\trCauchy{}公式的一个初等推论：设 $\bar{U}=\mathrm{int}(D)^q\times\RR^p$，$U=(\mathrm{int}(D)\setminus\{0\})^q\times\RR^p$，而对 $r\in\RR^p$ 令 $\bar{U}_r=\mathrm{int}(D)^q\times\{r\}$。设有函数 $g:\bar{U}\to\CC$，它在 $U$ 上\trSmooth，并且它到每个 $\bar{U}_r$ 的限制都是\trHolomorphic{}的。那么 $g$ 是\trSmooth{}的。最后，为证明 $\bar\gamma^T$ 是\trLocalDiffeomorphism，我们需要在所有这些参数皆为零的点处计算它关于诸\trGluingParameter{}的导数。在这里，复几何（更确切地说，带\trNode{}的\trRiemann\trSurface{}的\trDeformation{}理论）同样使它成为一道容易的练习。
\end{proof}
有了\cref{thm:dm-smooth-structure}在手，便容易看出，$\bar{\scr{R}}^{d+1}$ 的\trTopology{}与\trSmoothStructure——按其原初定义（对\trStriplikeEnd{}的任意\trUniversal{}选取，经由\trGluingMap{}\eqref{eq:extended-gluing-map} 给出）——与由\trEmbedding{}$\bar{\scr{R}}^{d+1}\subset\bar{\scr{M}}_{0,d+1}(\RR)$ 给出的那些相同。这就蕴含了上述全部性质。
定义 $\scr{S}^T\to\scr{R}^T$ 为诸\trUniversal{}\trFibration{}$\scr{S}^{|v|}\to\scr{R}^{|v|}$（$v\in\mathrm{Ve}(T)$）的\trPullback{}的\trDisjointUnion；$\scr{S}^T$ 的诸\trFiber{}将是 $|v|$ \trPointedDisc{}的\trDisjointUnion。\trUniversal\trFamily{}$\scr{S}^{d+1}\to\scr{R}^{d+1}$ 容许一个\trPartialCompactification，在集合论的层面上它就是
\begin{equation}\label{eq:universal-partial-compactification}
\bar{\scr{S}}^{d+1}=\coprod_T\scr{S}^T\longrightarrow\bar{\scr{R}}^{d+1}.
\end{equation}
在诸\trFamily{}的\trTotalSpace{}上使用\trGluingMap，便可以为它配备\trTopology{}与\trSmoothStructure。按定义，与\eqref{eq:gluing-family} 相关联的\trGluingMap{}$\gamma^T$ 被一个\trFibrewise\trIsomorphism{}所\trCovering
\begin{equation*}
\Gamma^T:\scr{S}\longrightarrow\scr{S}^{d+1}.
\end{equation*}
正如在解释一般\trGluingConstruction{}时所指出的，$\scr{S}$ 有一个自然的\trPartialCompactification{}$\bar{\scr{S}}\to\bar{\scr{R}}$。把它与\eqref{eq:universal-partial-compactification} 相比较即可看出，$\Gamma^T$ 可扩张为 $\bar\Gamma^T:\bar{\scr{S}}\to\bar{\scr{S}}^{d+1}$。与前面一样，由这些映射诱导的\trTopology{}与\trSmoothStructure，与把 $\bar{\scr{S}}^{d+1}$ 经\trEmbedding{}置入 $\bar{\scr{M}}_{0,d+1}$ 之上的\trStableCurve{}的\trUniversal\trFamily{}所得者相同。
